\chapter{Endpoint-Conditioned Movement Reconstruction}
\label{chap:interpolation}

\section{Intuition before formalism}

\begin{itemize}
  \item The model reconstructs uncertainty between two observed fixes rather than guessing one definitive missing route.
  \item Space is divided into grid cells and elapsed time into a finite number of movement transitions.
  \item A transition kernel specifies how the animal could move locally at each step.
  \item Conditioning retains only paths that start and end at the observed cells at the required times.
  \item Dynamic programming makes this conditioning feasible without enumerating every complete path \parencite{buchin2024randomwalks,hagedoornPhD}.
\end{itemize}

\section{Forward computation}

\begin{itemize}
  \item Initialize all probability mass at the observed start cell.
  \item At every model step, distribute mass according to the active Brownian, correlated, state-dependent, and terrain-adjusted transition matrix.
  \item Retain incoming heading as part of the state when correlated motion is used.
  \item Interpret mass at the observed endpoint after $T$ transitions as the probability of reaching that endpoint under the unconditioned local model.
  \item Treat zero endpoint mass as an unreachable configuration, not a numerical value that can safely be normalized.
\end{itemize}

\section{Backward sampling}

\begin{itemize}
  \item Begin at the observed endpoint.
  \item Select each predecessor with probability proportional to its forward mass and its probability of transitioning to the current state.
  \item Continue until the known start point is reached.
  \item Repetition yields representative samples from the discretized endpoint-conditioned model.
  \item Sampled walks represent uncertainty; they are not replicated observations or claims about the exact route.
\end{itemize}

\section{Connecting physical and discrete scales}

\begin{itemize}
  \item Derive the initial transition count from observed elapsed time and the native physical interval stored with the fitted kernel.
  \item Map the complete physical kernel support onto the realized grid instead of clipping it to a fixed number of cells.
  \item Current adaptive policy:
  \begin{itemize}
    \item $S=\lceil R/c\rceil$ for physical kernel radius $R$ and cell size $c$;
    \item use the elapsed-time transition count when the endpoint is reachable;
    \item increase $T$ only to the minimum required for geometric reachability when needed;
    \item record every adjustment because it changes the time interpretation.
  \end{itemize}
  \item Keep correlated kernels at their inferred native interval; allow explicit time transformation only for Brownian models.
\end{itemize}

\section{Failures and numerical safeguards}

\begin{itemize}
  \item Check finite-window boundaries, hard barriers, invalid start/end cells, absent state kernels, and inconsistent terrain mappings before computation.
  \item Use layer normalization, stored scale factors, or equivalent safeguards against probability underflow on long intervals.
  \item Require consistent semantics between forward densities, backward sampling, utilization distributions, and visit calculations.
  \item Return actionable diagnostics for no-path intervals instead of substituting a straight segment without disclosure.
  \item Verify serial, parallel, and GPU implementations against small exhaustively enumerable cases.
\end{itemize}

\section{Outputs}

\begin{itemize}
  \item Interpolated trajectories with timestamps, animal identifiers, walk-version identifiers, and segment boundaries.
  \item Optional serialized walks containing grid dimensions, terrain, steps, probability summaries, and metadata.
  \item Diagnostics for model interval, $T$, $S$, reachability, resampling, fallback behavior, and runtime.
  \item Static plots and interactive map layers that separate original fixes from reconstructed samples.
\end{itemize}
